\section{ماژول lillama\_compressor (LiLLaMA Compressor)}
\label{sec:module_lillama_compressor}

\subsection{معرفی و انگیزه}
این ماژول در چارچوب پژوهش حاضر پیاده‌سازی و ارزیابی شده و فاقد مقاله مستقل جداگانه در ادبیات تحقیق است.

تقریب ماتریس‌های تبدیل خطی و توجه در مدل‌های زبانی بزرگ از طریق تجزیه مقادیر منفرد (\lr{SVD}) در صورتی که صرفاً مبتنی بر ماتریس وزن ایستا $\mathbf{W}$ انجام گیرد، منجر به از دست رفتن جهت‌های بحرانی اطلاعات در فضای فعال‌سازی می‌گردد. دلیل بنیادین این پدیده آن است که مقادیر منفرد ماتریس وزن نشان‌دهنده میزان اهمیت مؤلفه‌ها در فضای پویا و داده‌محور فعال‌سازی‌های پنهان ورودی و خروجی نیستند. هرس یا تقریب زیرفضاهایی که فعال‌سازی‌های حجیم یا پرتاثیر را جابجا می‌کنند، به رانش بازنمایی شدید (\lr{Representation Drift}) و تخریب یکپارچگی توجه منجر می‌شود.

فشرده‌ساز \lr{LiLLaMA (Local Distillation via Activation Cross-Covariance SVD)} رویکردی داده‌محور و کالیبره‌شده برای حل این چالش ارائه می‌دهد. ایده بنیادین در این ماژول، تقطیر محلی لایه‌ها بر پایه ماتریس هم‌کوواریانس (\lr{Cross-Covariance Matrix}) میان تانسور فعال‌سازی ورودی و خروجی لایه است. با اعمال تجزیه مقادیر منفرد بر این تانسور هم‌کوواریانس، پروجکشن‌های خطی جدیدی بازسازی می‌شوند که کمترین خطای تجربی را در بازتولید بردار خروجی به ازای ورودی‌های کالیبراسیون ایجاد نموده و پایداری شبکه را در نسبت‌های فشرده‌سازی بالا حفظ می‌کنند.

\subsection{مبانی نظری و فرمول‌بندی ریاضی}
فرض کنید برای یک لایه خطی با ماتریس وزن $\mathbf{W} \in \mathbb{R}^{d_{\text{out}} \times d_{\text{in}}}$، جفت فعال‌سازی‌های ورودی و خروجی در حین اجرای گذر پیشرو بر روی مجموعه داده کالیبراسیون به ترتیب $\mathbf{X} \in \mathbb{R}^{N \times d_{\text{in}}}$ و $\mathbf{Y} \in \mathbb{R}^{N \times d_{\text{out}}}$ باشند.

ماتریس هم‌کوواریانس تجربی متقابل میان خروجی و ورودی لایه به صورت زیر محاسبه می‌شود:
\begin{equation}
\mathbf{C}_{YX} = \frac{1}{N} \mathbf{Y}^\top \mathbf{X} \in \mathbb{R}^{d_{\text{out}} \times d_{\text{in}}}
\end{equation}

این ماتریس رابطه ساختاری میان ویژگی‌های فضای ورودی و خروجی تحت توزیع داده‌های واقعی را کپسوله‌سازی می‌کند. سپس بر روی ماتریس $\mathbf{C}_{YX}$، تجزیه مقادیر منفرد کوتاه و بریده‌شده (\lr{Truncated SVD}) اعمال می‌گردد:
\begin{equation}
\mathbf{C}_{YX} \approx \mathbf{U}_r \mathbf{\Sigma}_r \mathbf{V}_r^\top
\end{equation}
که در آن $\mathbf{U}_r \in \mathbb{R}^{d_{\text{out}} \times r}$، $\mathbf{\Sigma}_r \in \mathbb{R}^{r \times r}$ و $\mathbf{V}_r \in \mathbb{R}^{d_{\text{in}} \times r}$ بوده و رتبه موثر $r$ با تحلیل نسبت انرژی تجمعی مقادیر منفرد تعیین می‌شود:
\begin{equation}
E(r) = \frac{\sum_{i=1}^r \sigma_i^2}{\sum_{j=1}^{\min(d_{\text{in}}, d_{\text{out}})} \sigma_j^2} \ge 0.95
\end{equation}

وزن بازسازی‌شده تقطیریافته $\mathbf{W}_{\text{distill}}$ به صورت زیر حاصل می‌شود:
\begin{equation}
\mathbf{W}_{\text{distill}} = \mathbf{U}_r \mathbf{\Sigma}_r \mathbf{V}_r^\top
\end{equation}

خطای فروبنیوس نسبی تقریب میان وزن اصلی و وزن تقطیریافته محاسبه می‌گردد:
\begin{equation}
\mathcal{E}_{\text{rel}} = \frac{\|\mathbf{W} - \mathbf{W}_{\text{distill}}\|_F}{\|\mathbf{W}\|_F + \epsilon}
\end{equation}
چنانچه $\mathcal{E}_{\text{rel}} > \text{max\_error}$ باشد (با مقدار پیش‌فرض $0.20$)، ماژول به صورت خودکار به وزن اصلی لایه رجوع کرده (\lr{Fallback to Original}) تا از تخریب ساختاری مدل جلوگیری به عمل آید.

\subsection{معماری و پیاده‌سازی}
کلاس \lr{LiLLaMACompressor} در مسیر \lr{Codes/framework/compression\_engine/lillama\_compressor.py} واقع شده و تحت شناسه‌های \lr{lillama}، \lr{lillama\_compressor}، \lr{lillama\_class} و \lr{lillama\_factory} در رجیستری فشرده‌سازی ثبت شده است. مؤلفه‌ها و مکانیزم‌های کلیدی پیاده‌سازی عبارتند از:
\begin{itemize}
    \item \textbf{مکانیزم قلاب‌های رهگیری پویا (\lr{Forward Activation Hooks}):} ماژول قلاب‌های موقتی را بر روی تمامی لایه‌های خطی هدف نصب نموده و در یک گذر پیشرو بر روی $N_{\text{batch}}$ دسته داده متنی، تانسورهای فعال‌سازی ورودی و خروجی را جمع‌آوری و برای مدیریت حافظه به پردازنده میزبان انتقال می‌دهد.
    \item \textbf{سازگاری کامل با لایه‌های \lr{Conv1D} و شبکه‌های دیکودر:} لایه‌های پیچشی ۱بعدی (نظیر ساختار \lr{c\_attn} در \lr{GPT-2}) از طریق ترانهاده‌سازی خودکار و بازآرایی ابعاد پشتیبانی می‌شوند.
    \item \textbf{فیلتر هوشمند لایه‌ها و لایه‌های بحرانی:} لایه‌های حساس نظیر سر رده‌بندی، لایه‌های نشاندن کلمات (\lr{Embeddings}) و لایه‌های با ابعاد کمتر از $\text{min\_layer\_dim}=8$ از فرآیند فشرده‌سازی مستثنی می‌شوند تا پایداری خروجی تضمین گردد.
    \item \textbf{تضمین عدم اجرای تکراری (\lr{Idempotency Guard}):} با تنظیم پرچم اختصاصی بر روی مدل پس از فشرده‌سازی، از اجرای ناخواسته فشرده‌سازی چندگانه جلوگیری می‌شود.
\end{itemize}

\subsection{نتایج تجربی}
ارزیابی ماژول \lr{lillama\_compressor} بر روی مدل‌های زبانی پایه در چارچوب بنچ‌مارک استاندارد و مقادیر ثبت‌شده در \lr{thesis\_final/master\_verified\_results.json} در جدول \ref{tab:benchmark_lillama_compressor} ارائه شده است:

\begin{table}[H]
\centering
\small
\caption{نتایج تجربی و اعتبارسنجی بنچ‌مارک فشرده‌ساز \lr{lillama\_compressor} روی مدل‌های پایه}
\label{tab:benchmark_lillama_compressor}
\begin{tabular}{lccccc}
\toprule
\textbf{مدل پایه} \lr{(Base Model)} & \textbf{دقت} \lr{(Accuracy)} & \textbf{شاخص} \lr{(F1-Score)} & \textbf{حافظه} \lr{(Peak VRAM)} & \textbf{توان} \lr{(Throughput)} & \textbf{وضعیت} \\
\midrule
\lr{DistilBERT} & \lr{87.50\%} & \lr{0.8868} & \lr{337.59 MB} & \lr{187,624.63 tok/s} & \lr{PASS} \\
\lr{RoBERTa} & \lr{85.44\%} & \lr{0.8603} & \lr{583.49 MB} & \lr{100,764.43 tok/s} & \lr{PASS} \\
\lr{GPT-2} & \lr{56.88\%} & \lr{0.6835} & \lr{634.91 MB} & \lr{70,009.59 tok/s} & \lr{PASS} \\
\lr{Qwen} & \lr{49.08\%} & \lr{---} & \lr{---} & \lr{---} & \lr{PASS} \\
\bottomrule
\end{tabular}
\end{table}

داده‌های تجربی نشان می‌دهند که ماژول بر روی مدل‌های انکودر به سطح بسیار بالایی از کارایی دست یافته است؛ به طوری که در \lr{DistilBERT} دقت \lr{87.50\%} همراه با شاخص F1 برابر با \lr{0.8868} و توان پردازشی بالا \lr{187,624.63 tok/s} را به ثبت رسانده است. در مدل \lr{RoBERTa} نیز دقت \lr{85.44\%} با پایداری کامل حاصل شده است. با این حال، در مدل‌های اتورگرسیو نظیر \lr{GPT-2} (\lr{56.88\%}) و \lr{Qwen} (\lr{49.08\%}) افت دقت مشاهده می‌شود.

\subsection{محدودیت‌ها و تحلیل شکست}
تحلیل مکانیزم عملکردی و خطاهای تجربی ماژول نکات زیر را آشکار می‌سازد:
\begin{itemize}
    \item \textbf{حساسیت بالا به پوشش مجموعه داده کالیبراسیون:} تقطیر وزن‌ها بر پایه ماتریس هم‌کوواریانس تجربی $\mathbf{C}_{YX}$ وابستگی مستقیمی به توزیع آماری داده‌های ورودی کالیبراسیون دارد؛ در صورت محدود بودن دسته‌ها، ماتریس هم‌کوواریانس دچار سوگیری شده و توان تعمیم مدل کاهش می‌یابد.
    \item \textbf{انباشت خطا در فرآیند تولید اتورگرسیو:} در مدل‌های دیکودر لایه‌به‌لایه‌محور (\lr{GPT-2} و \lr{Qwen})، تقریب وزن‌ها با بازنویسی مستقیم وزن بدون فاکتورگیری تفکیک‌شده، خطاهای کوچکی در گام‌های اولیه بازتولید می‌کند که در گذر گام‌های متوالی تولید بازگشتی بزرگ‌نمایی شده و افت دقت را تشدید می‌نماید.
    \item \textbf{مصرف حافظه حین فاز جمع‌آوری فعال‌سازی‌ها:} تجمیع تانسورهای ورودی و خروجی تمام لایه‌ها در رم سیستم نیازمند تخصیص پویای منابع و تخلیه مکرر بافرها در مدل‌های بسیار عمیق است.
\end{itemize}
